\needspace{10\baselineskip}
\hypertarget{native-317-zusatz15}{}
\begin{center}
\pdfbookmark[1]{15. Differentialformeln bei Berechnung der Oppositionen}{zusatz-015}
[15.]
\end{center}
\begin{center}
Differentialformeln bei Berechnung der Oppositionen.
\end{center}
\begin{center}
\begin{tabular}{@{}ll@{}}
\(\lambda\) & ber. hel. Länge\\
\(\lambda^{\prime}\) & beob. hel. Länge\\
\(\omega\) & Sonnennähe\\
\(v\) & wahre Anomalie\\
\(E\) & excentrische Anomalie\\
\(e=\sin\varphi\) & Excentricität\\[5pt]
\(r\) & Radius Vector\\
\(a\) & halbe grosse Axe\\
\(i\) & Neigung der Bahn\\
\(\AscendingNode\) & Knoten\\
\(\vartheta\) & heliocentrische Breite\\
\(\beta\) & geocentrische Breite\\[5pt]
\(n\) & Tage seit Epoche\\
\(\mathord{7}\) & tägliche Bewegung
\end{tabular}
\end{center}
\[
\begin{aligned}
\frac{\sin2(\lambda-\AscendingNode)}{\sin2(v+\omega-\AscendingNode)}&=A,\\[5pt]
\frac{Aaa\cos\varphi}{rr}&=B,\\[5pt]
\frac{Aaa\sin E(2-ee-e\cos E)}{rr}&=C,
\end{aligned}
\]
\[
\begin{aligned}
\mathrm{d}\lambda={}&B\mathord{.}\mathrm{dEp.}+nB\mathord{.}\mathrm{d}\mathord{7}+(A-B)\mathrm{d}\omega+C\mathord{.}\mathrm{d}\varphi\\
&+(1-A)\mathrm{d}\AscendingNode-\tfrac12\operatorname{tang}i\sin2(\lambda-\AscendingNode)\mathrm{d}i;
\end{aligned}
\]
\[
\frac{\sin\beta\sin(\beta-\vartheta)}{\sin\vartheta}=D,\qquad
\sin\beta\cos(\beta-\vartheta)\cos\vartheta=\mathrm{E},
\]
\[
\frac{Da\operatorname{tang}\varphi\sin v}{r}=F,\qquad
\frac{Da\cos\varphi\cos v}{r}=G,
\]
\[
\begin{aligned}
\mathrm{d}\beta={}&-F\mathord{.}\mathrm{dEp.}+\left(\frac{2D}{3\mathord{7}}-nF\right)\mathrm{d}\mathord{7}+F\mathord{.}\mathrm{d}\omega+G\mathord{.}\mathrm{d}\varphi\\
&-\frac{\mathrm{E}}{\operatorname{tang}(\lambda^{\prime}-\AscendingNode)}\mathrm{d}\AscendingNode
+\frac{\mathrm{E}}{\sin i\cos i}\mathrm{d}i.
\end{aligned}
\]
\SourceLeaf{318}
\begin{center}
Beispiel. Opposition der Pallas von 1807.
\end{center}
[Mit den in der Disquisitio de elementis ellipticis Palladis art. 2. 3. 8 gegebenen Werthen
\begin{center}
\begin{tabular}{@{}ll@{}}
\(\log a=0{,}44223\) & \([v=107^{0}\,47^{\prime}\,31^{\prime\prime}\)\\
\(e=0{,}24476\) & \(\lambda=223\quad37\quad25\)\\
\(\varphi=14^{0}\,10^{\prime}\,4^{\prime\prime}\) & \(\lambda^{\prime}=223\quad37\quad28]\)\\
\(\tilde\omega=121\quad8\quad59\) & \(\beta=42\quad11\quad28\)\\
\(\AscendingNode=172\quad32\quad24\) & \(\vartheta=28\quad14\quad51\)\\
\(i=34\quad37\quad32]\) &
\end{tabular}
\end{center}
[und den leicht zu findenden
\[
E=93^{0}\,46^{\prime}\,4^{\prime\prime}\qquad \log r=0{,}44916
\]
ergibt die Rechnung:]
\begin{center}
\renewcommand{\arraystretch}{1.4}
\begin{tabular}{@{}lr@{}}
\(\sin2(\lambda-\AscendingNode)\) . . & \(9{,}99013\)\\
\(\sin2(v+\omega-\AscendingNode)\)  & \(9{,}96466\)\\
\(-\tfrac12\operatorname{tang}i\) . . . & \(9{,}53814_n\)\\
\hline
\(A\) . . . . . . . . . . & \(0{,}02547\)\\
\(\dfrac{rr}{aa}\) . . . . . . . . . & \(0{,}01386\)\\[4pt]
\(\dfrac{Aaa}{rr}\) . . . . . . . . & \(0{,}01161\)\\
\(\sin E\) . . . . . . . & \(9{,}99906\)\\
\(2-ee-e\cos E\)  & \(0{,}29141\)\\
\hline
\(C\) . . . . . . . . . . & \(0{,}30208\)\\
\(B\) . . . . . . . . . . & \(9{,}99819\)\\
\(n\) . . . . . . . . . . & \(3{,}20020\)\\
\end{tabular}
\end{center}
\begin{center}
\renewcommand{\arraystretch}{1.4}
\begin{tabular}{@{}lr@{}}
\(\sin\beta\) . . . . . . & \(9{,}82712\)\\
\(\sin(\beta-\vartheta)\) . . & \(9{,}38195\)\\
\(\mathrm{C}\log\sin\vartheta\) . & \(0{,}32488\)\\
\hline
\(D\) . . . . . . . . & \(9{,}53395\)\\
\(\dfrac{2}{3\mathord{7}}\) . . . . . . . . & \(2{,}25167\)\\[4pt]
\(\dfrac{r}{a}\) . . . . . . . . & \(0{,}00693\)\\[4pt]
\(\dfrac{Da}{r}\) . . . . . . . & \(9{,}52702\)\\
\(\operatorname{tang}\varphi\) . . . . . & \(9{,}40216\)\\
\(\sin v\) . . . . . . & \(9{,}97872\)\\
\hline
\(F\) . . . . . . . . & \(8{,}90790\)\\
\(n\) , . . . . . . . . & \(3{,}20020\)\\
\end{tabular}
\end{center}
\begin{center}
\renewcommand{\arraystretch}{1.4}
\begin{tabular}{@{}lr@{}}
\(\sin\beta\) . . . . . . & \(9{,}82712\)\\
\(\cos(\beta-\vartheta)\) . . & \(9{,}98701\)\\
\(\cos\vartheta\) . . . . . . & \(9{,}94493\)\\
\hline
\(\mathrm{E}\) . . . . . . . . & \(9{,}75906\)\\
\(\operatorname{tang}(\lambda^{\prime}-\AscendingNode)\)  & \(0{,}09293\)\\
\(\tfrac12\sin2i\) . . . . & \(9{,}66985\)\\
\(\dfrac{Da}{r}\) . . . . . . . & \(9{,}52702\)\\
\([\cos\varphi\) . . . . . & \(9{,}98658]\)\\
\(\cos v\) . . . . . . & \(9{,}48510_n\)\\
\hline
\([G\) . . . . . . . & \(8{,}99870_n]\)\\
\end{tabular}
\end{center}
\[
\begin{aligned}
\mathrm{d}\lambda={}&0{,}99584\,\mathrm{dEp.}+1579{,}03\,\mathrm{d}\mathord{7}+2{,}00486\,\mathrm{d}\varphi\\
&+0{,}06456\,\mathrm{d}\omega-0{,}06040\,\mathrm{d}\AscendingNode-0{,}33750\,\mathrm{d}i
\end{aligned}
\]
\[
\begin{aligned}
\mathrm{d}\beta={}&-0{,}08089\,\mathrm{dEp.}-67{,}22\,\mathrm{d}\mathord{7}-0{,}09970\,\mathrm{d}\varphi\\
&+0{,}08089\,\mathrm{d}\omega-0{,}46359\,\mathrm{d}\AscendingNode+1{,}22803\,\mathrm{d}i.
\end{aligned}
\]
\par\medskip\begin{center}\rule{.22\textwidth}{.3pt}\end{center}
